\section{Introduction}

Bipolar Disorder affects roughly 1--2\% of the global population, yet studies consistently find that patients wait an average of nearly a decade before receiving a correct diagnosis \cite{garcia2018depresjon}. The core reason for this delay is structural: people experiencing a depressive episode are the ones who seek clinical help, while the hypomanic or manic highs that define BD often feel productive or even pleasant, and so go unreported. A clinician who sees only the depressive presentation has very little to go on when trying to decide whether they are looking at Bipolar Disorder or ordinary Major Depressive Disorder (MDD). This is not just a diagnostic inconvenience it is a clinical hazard. Standard antidepressant therapy prescribed to a patient with unrecognized BD can trigger manic switching, rapid cycling, or mixed states that are considerably harder to manage. There is therefore a genuine and urgent need for an objective, passive tool that can flag the physiological signature of bipolar depression before the wrong treatment is started.

Current diagnostic practice relies almost entirely on structured clinical interviews and patient-reported questionnaires such as the PHQ-9, the MADRS, or the HCL-32. These instruments have real value, but they share a common weakness: they depend on the patient's ability and willingness to accurately recall their own mood states over the preceding days or weeks. Recall bias is well documented in this population, and the very nature of a depressive episode low energy, reduced motivation, cognitive slowing makes retrospective self-report unreliable. On the computational side, a handful of studies have applied classical machine learning methods like Support Vector Machines and Random Forests to actigraphy data \cite{zanella2019feature}, and these represent a meaningful step forward. However, they all rely on manually extracted statistical summaries mean activity, variance, zero-crossing rate computed over fixed windows. These hand-crafted features can capture the average level of motor activity, but they tend to miss the subtler temporal structure: the fragmented sleep bouts, the irregular rest-activity transitions, and the multi-day rhythmic drift that may differ systematically between bipolar and unipolar depression.

This project asks a straightforward but underexplored question: can a deep sequence model learn to make the bipolar-versus-unipolar distinction directly from raw actigraphy, without any manual feature engineering? We use the publicly available Depresjon dataset \cite{garcia2018depresjon}, which provides continuous wrist actigraphy recorded at one-minute resolution across 23 patients with clinically confirmed mood disorder diagnoses and 32 healthy controls. Our proposed architecture pairs a 1D Convolutional Neural Network, which is well-suited to picking up localized kinetic patterns within short temporal windows, with a Long Short-Term Memory network that can integrate information across the full multi-day recording. The combination allows the model to simultaneously attend to micro-scale movement bursts the kind of restless, low-amplitude activity that disrupts sleep and to macro-scale circadian shifts that play out over days. We believe this design better matches the actual temporal structure of the clinical signal than any fixed-window statistical approach, and we run two controlled experiments to test that belief.

\section*{}
